\documentclass[11pt]{article}

\usepackage[margin=1in]{geometry}
\usepackage{amsmath,amssymb,amsthm,mathtools}
\usepackage{booktabs,longtable,tabularx,array}
\usepackage{enumitem}
\usepackage{xcolor}
\usepackage{hyperref}
\usepackage{fancyhdr}
\usepackage{microtype}
\usepackage[T1]{fontenc}
\usepackage{lmodern}

\hypersetup{
  colorlinks=true,
  linkcolor=blue!55!black,
  urlcolor=blue!55!black,
  citecolor=blue!55!black,
  pdftitle={Aut(F2) and ACA Equivalence: A Formal Tutorial and Practice Workbook},
  pdfauthor={ACSolverX equivalence-class analysis}
}

\pagestyle{fancy}
\fancyhf{}
\lhead{\(\operatorname{Aut}(F_2)\) and ACA Equivalence}
\rhead{\thepage}
\setlength{\headheight}{14pt}
\setlength{\parindent}{0pt}
\setlength{\parskip}{0.55em}
\setlist{itemsep=0.3em,topsep=0.4em}

\newtheorem{theorem}{Theorem}[section]
\newtheorem{proposition}[theorem]{Proposition}
\newtheorem{lemma}[theorem]{Lemma}
\newtheorem{corollary}[theorem]{Corollary}
\theoremstyle{definition}
\newtheorem{definition}[theorem]{Definition}
\newtheorem{example}[theorem]{Example}
\newtheorem{exercise}{Exercise}[section]
\theoremstyle{remark}
\newtheorem{remark}[theorem]{Remark}
\newtheorem{warning}[theorem]{Warning}

\newcommand{\F}{F_2}
\newcommand{\Aut}{\operatorname{Aut}}
\newcommand{\End}{\operatorname{End}}
\newcommand{\Ncl}[1]{\left\langle\!\left\langle #1\right\rangle\!\right\rangle}
\newcommand{\AC}{\mathrm{AC}}
\newcommand{\ACA}{\mathrm{ACA}}
\newcommand{\id}{\operatorname{id}}
\newcommand{\rot}{\operatorname{rot}}
\newcommand{\canon}{\operatorname{canon}}
\newcommand{\len}[1]{\left|#1\right|}
\newcommand{\eps}{\varepsilon}
\newcommand{\sourcefile}[1]{\path{#1}}
\newcommand{\diagnostic}[1]{\par\smallskip\textit{Diagnostic: #1}}

\begin{document}
\hypersetup{pageanchor=false}
\pagenumbering{roman}

\begin{titlepage}
\centering
\vspace*{1.3in}
{\Huge\bfseries \(\operatorname{Aut}(F_2)\) and ACA Equivalence\par}
\vspace{0.35in}
{\Large A Formal Tutorial and Practice Workbook\par}
\vspace{0.2in}
{\large Why 261 benchmark presentations reduce to at most 124 distinct problems\par}
\vfill
{\large ACSolverX equivalence-class analysis\par}
\vspace{0.15in}
{\large July 2026\par}
\vfill
\begin{minipage}{0.86\textwidth}
\small
\textbf{Central claim.}
The computation does not prove that any of the 261 presentations is
AC-trivial.  It proves that a finite graph of certified Andrews--Curtis and
change-of-basis edges on the 261 inputs has 124 connected components.
Consequently, proving one representative from each component AC-trivial
would prove all 261 AC-trivial.  Each verified component is contained in a
true ACA class, so 124 is an upper bound on the number of distinct problems,
not a proof that 124 mutually inequivalent AC-classes exist.
\end{minipage}
\end{titlepage}

\hypersetup{pageanchor=true}
\tableofcontents
\clearpage
\pagenumbering{arabic}

\section*{How to use this guide}
\addcontentsline{toc}{section}{How to use this guide}

This document has three simultaneous purposes.

\begin{enumerate}
  \item It develops the group theory needed to understand why
  \(\Aut(\F)\) is relevant to balanced presentations.
  \item It proves, move by move, why a change of free basis preserves
  AC-triviality and gives quantitative upper bounds on the transported
  number of elementary AC moves.
  \item It explains the complete computational pipeline
  \[
    261\longrightarrow261\longrightarrow171\longrightarrow168
    \longrightarrow126\longrightarrow125\longrightarrow124
  \]
  and separates exact conclusions from search-dependent upper bounds.
\end{enumerate}

Every substantive section ends with exercises.  A complete answer key,
including diagnostic advice, begins in Section~\ref{sec:answers}.  The final
summary in Section~\ref{sec:presentation} is designed for a lab presentation.

\begin{warning}[No silent theorem dependencies]
Elementary free-group and AC facts are proved here.  Whitehead peak reduction
and the Nielsen basis theorem are classical results too large to reprove
honestly in a tutorial of this scope.  They are stated with their hypotheses,
their precise role is isolated, and every project-specific conclusion drawn
from them is proved.
\end{warning}

\section{Conventions: three elementary AC moves}
\label{sec:conventions}

There are multiple numbering conventions in the literature and in the
ACSolverX materials.  Mixing them changes move counts.  This guide uses the
convention of the source paper
\sourcefile{literature/txt/math_ml_paper_2408.15332.txt}, lines 296--298.

\begin{definition}[Source-paper elementary AC moves]
Let
\[
  P=\langle x_1,\ldots,x_n\mid r_1,\ldots,r_n\rangle
\]
be balanced.  For \(i\ne j\), the elementary moves are:
\begin{align*}
  \text{\rm(AC1)}&\quad r_i\longmapsto r_i r_j,\\
  \text{\rm(AC2)}&\quad r_i\longmapsto r_i^{-1},\\
  \text{\rm(AC3)}&\quad r_i\longmapsto g r_i g^{-1},
       \qquad g\in\{x_1^{\pm1},\ldots,x_n^{\pm1}\}.
\end{align*}
Only one relator is replaced; all other relators remain unchanged.
\end{definition}

\begin{center}
\begin{tabularx}{\textwidth}{>{\bfseries}l X X}
\toprule
Operation & Source-paper convention used here & Two-Hump/solver prose convention\\
\midrule
Combine & AC1: \(r_i\mapsto r_i r_j\) & Often called AC2\\
Invert & AC2: \(r_i\mapsto r_i^{-1}\) & Often called AC1\\
Conjugate & AC3 by one generator or inverse & Often AC3 by an arbitrary word\\
Solver S-move & A finite sequence of AC1--AC3 & One encoded macro-move\\
\bottomrule
\end{tabularx}
\end{center}

\begin{warning}[No inverse multiplier in elementary AC1]
The source-paper AC1 is \(r_i\mapsto r_i r_j\), not
\(r_i\mapsto r_i r_j^{-1}\).  The latter can be realized while restoring
\(r_j\) by three moves:
\[
(r_i,r_j)\xrightarrow{\AC2}(r_i,r_j^{-1})
\xrightarrow{\AC1}(r_i r_j^{-1},r_j^{-1})
\xrightarrow{\AC2}(r_i r_j^{-1},r_j).
\]
Thus the solver parameter \(j=-1\) is part of a macro-move under the
source-paper convention.
\end{warning}

Relator order is not mathematical data, so swapping \(r_1,r_2\) is treated
as zero-cost bookkeeping.  Free reduction is also zero-cost because it
replaces a word by the same element of the free group.

\subsection*{Exercises}

\begin{exercise}
Using the convention above, write a three-move sequence that reverses one
positive AC1 move:
\[
  (r_1r_2,r_2)\longrightarrow (r_1,r_2).
\]
\end{exercise}

\begin{exercise}
Explain why the phrase ``one rotation move'' is ambiguous unless one states
whether AC3 allows conjugation by an arbitrary word or only by one generator.
\end{exercise}

\begin{exercise}
Classify each operation as elementary, macro, or bookkeeping:
\[
r_1\mapsto r_1^{-1},\qquad
r_1\mapsto r_1r_2^{-1},\qquad
r_1\mapsto w r_1w^{-1},\qquad
(r_1,r_2)\mapsto(r_2,r_1).
\]
Use the source-paper convention.
\end{exercise}

\section{Free groups and balanced presentations}
\label{sec:free-presentations}

\subsection{The free group \(F_2\)}

\begin{definition}[Free group on two generators]
The free group
\[
  \F=\langle x,y\rangle
\]
consists of reduced words in the alphabet
\[
  \{x,X,y,Y\},\qquad X=x^{-1},\quad Y=y^{-1}.
\]
A word is reduced when it contains no adjacent inverse pair
\(xX,Xx,yY,Yy\).  Multiplication is concatenation followed by free
reduction.
\end{definition}

For a word \(w=a_1\cdots a_m\), its inverse is
\[
  w^{-1}=a_m^{-1}\cdots a_1^{-1}.
\]
Thus, for example,
\[
  (xyX)^{-1}=xYX.
\]
The empty word is the identity.  We write \(\len{w}\) for the length of the
freely reduced word representing \(w\).

\begin{theorem}[Universal property of \(\F\)]
\label{thm:universal}
Let \(H\) be any group and choose \(a,b\in H\).  There exists a unique
homomorphism \(f:\F\to H\) satisfying \(f(x)=a\) and \(f(y)=b\).
\end{theorem}

\begin{proof}
Define \(f\) on a word by replacing each letter:
\[
x\mapsto a,\quad X\mapsto a^{-1},\quad
y\mapsto b,\quad Y\mapsto b^{-1},
\]
and multiplying the resulting elements of \(H\).  Deleting an adjacent
inverse pair in \(\F\) deletes \(aa^{-1}\), \(a^{-1}a\), \(bb^{-1}\), or
\(b^{-1}b\) in \(H\), so the value is unchanged by free reduction.
Therefore the map is well-defined.  Concatenation gives
\(f(uv)=f(u)f(v)\), so it is a homomorphism.  Any homomorphism with the
specified values on \(x,y\) must evaluate every word by the same rule,
which proves uniqueness.
\end{proof}

This theorem explains why assigning words to \(x\) and \(y\) automatically
defines an endomorphism \(\F\to\F\).  It does \emph{not} automatically define
an automorphism: invertibility must still be proved.

\subsection{Presentations are quotients of the free group}

\begin{definition}[Normal closure]
For \(r_1,r_2\in\F\), the normal closure
\[
  \Ncl{r_1,r_2}
\]
is the smallest normal subgroup of \(\F\) containing \(r_1,r_2\).
Equivalently, its elements are finite products of conjugates of
\(r_1^{\pm1}\) and \(r_2^{\pm1}\).
\end{definition}

\begin{definition}[Balanced two-generator presentation]
\[
  P=\langle x,y\mid r_1,r_2\rangle
\]
denotes the quotient group
\[
  G_P=\F/\Ncl{r_1,r_2}.
\]
It is balanced because the number of relators equals the number of
generators.  Balanced does not mean trivial.
\end{definition}

The presentation defines the trivial group precisely when
\[
  G_P=1
  \quad\Longleftrightarrow\quad
  \Ncl{r_1,r_2}=\F.
\]
The standard trivial presentation is
\[
  T=\langle x,y\mid x,y\rangle.
\]

\begin{definition}[AC-trivial]
A balanced presentation is AC-trivial when a finite sequence of elementary
AC moves carries its relator pair to the standard pair \((x,y)\), up to
free reduction and relator order.
\end{definition}

\begin{warning}[Trivial group versus AC-trivial presentation]
If a presentation is AC-trivial, it certainly presents the trivial group.
The converse is the Andrews--Curtis conjecture and is not known in general.
The 261 benchmark inputs present the trivial group, but none has been proved
AC-trivial by the equivalence-class search.
\end{warning}

\subsection*{Exercises}

\begin{exercise}
Freely reduce \(xyYXxXY\), and compute the inverse of the reduced word.
\end{exercise}

\begin{exercise}
Prove directly that
\[
  \langle x,y\mid xy,y\rangle
\]
presents the trivial group.  Then give a source-convention AC sequence to
\((x,y)\).
\end{exercise}

\begin{exercise}
Why does the equality \(G_P=1\) not by itself provide an AC path from
\((r_1,r_2)\) to \((x,y)\)?
\end{exercise}

\begin{exercise}
State exactly where balancedness is used in the definition of the
Andrews--Curtis problem, and explain why a two-generator presentation with
three relators is not one of the objects counted in the 261-to-124 pipeline.
\end{exercise}

\section{Automorphisms of \(F_2\) and their action on presentations}
\label{sec:aut}

\begin{definition}
An endomorphism \(\phi:\F\to\F\) is an automorphism if it is bijective.
The group of all such maps under composition is \(\Aut(\F)\).
\end{definition}

By Theorem~\ref{thm:universal}, words \(u,v\in\F\) determine a unique
endomorphism
\[
  \phi(x)=u,\qquad \phi(y)=v.
\]
To prove that \(\phi\) is an automorphism, it is enough to construct an
endomorphism \(\psi\) satisfying both
\[
  \psi\circ\phi=\id_{\F},
  \qquad
  \phi\circ\psi=\id_{\F}.
\]
Because homomorphisms out of \(\F\) are determined by \(x,y\), each equality
needs to be checked only on those two generators.

\begin{example}[The basic length-changing automorphism]
\label{ex:basic-aut}
Let
\[
  \phi(x)=xy,\qquad \phi(y)=y.
\]
The universal property makes this a homomorphism.  Define
\[
  \psi(x)=xY,\qquad \psi(y)=y.
\]
Then
\begin{align*}
\psi(\phi(x))&=\psi(xy)=(xY)y=x,&
\psi(\phi(y))&=y,\\
\phi(\psi(x))&=\phi(xY)=(xy)Y=x,&
\phi(\psi(y))&=y.
\end{align*}
Hence \(\psi=\phi^{-1}\), so \(\phi\in\Aut(\F)\).
\end{example}

\begin{remark}
The rule on inverse letters is forced:
\[
  \phi(X)=\phi(x^{-1})=\phi(x)^{-1},\qquad
  \phi(Y)=\phi(y)^{-1}.
\]
For Example~\ref{ex:basic-aut},
\[
  \phi(X)=(xy)^{-1}=YX,\qquad \phi(Y)=Y.
\]
One must reverse the image word and invert every letter.
\end{remark}

\subsection{Why the presented group is unchanged}

For \(P=\langle x,y\mid r_1,r_2\rangle\), define
\[
  \phi(P)=\langle x,y\mid\phi(r_1),\phi(r_2)\rangle.
\]
The symbols \(x,y\) remain the displayed free generators; the relator words
have been rewritten in a new free basis.

\begin{theorem}[Change of free basis induces a quotient isomorphism]
\label{thm:quotient-iso}
Let \(N=\Ncl{r_1,r_2}\).  If \(\phi\in\Aut(\F)\), then
\[
  \phi(N)=\Ncl{\phi(r_1),\phi(r_2)}
\]
and
\[
  \overline{\phi}:\F/N\longrightarrow \F/\phi(N),
  \qquad
  wN\longmapsto\phi(w)\phi(N)
\]
is an isomorphism.
\end{theorem}

\begin{proof}
An automorphism sends products, inverses, and conjugates to products,
inverses, and conjugates.  Therefore \(\phi(N)\) is a normal subgroup
containing \(\phi(r_1),\phi(r_2)\), so
\(\Ncl{\phi(r_1),\phi(r_2)}\subseteq\phi(N)\).
Applying the same argument to \(\phi^{-1}\) gives the reverse inclusion.

If \(wN=w'N\), then \(w^{-1}w'\in N\), hence
\(\phi(w)^{-1}\phi(w')\in\phi(N)\), so
\(\overline{\phi}\) is well-defined.  It is a homomorphism because \(\phi\)
is.  The map induced by \(\phi^{-1}\) is its two-sided inverse.
\end{proof}

\begin{corollary}
If \(P\) presents the trivial group, then \(\phi(P)\) presents the trivial
group.  The transformed presentation remains balanced.
\end{corollary}

\begin{warning}
Theorem~\ref{thm:quotient-iso} proves that \(P\) and \(\phi(P)\) present
isomorphic groups.  That alone does not prove either presentation
AC-trivial.  The stronger AC-triviality statement requires the transported
path theorem in Section~\ref{sec:transport}.
\end{warning}

\subsection{How the project verifies an automorphism}

The equivalence pipeline does not trust an arbitrary dictionary
\(\{x\mapsto u,y\mapsto v\}\).  It independently checks that \((u,v)\) is a
free basis using Nielsen reduction.  The classical Nielsen basis theorem
says that a generating pair is a basis exactly when elementary Nielsen
transformations reduce it to a signed permutation of \((x,y)\).
The verifier in
\sourcefile{experiments/equivalence_classes/verify/verify_certificates.py}
implements this check independently of the Whitehead code that produced the
witness.

\subsection*{Exercises}

\begin{exercise}
Prove that
\[
  \phi_k(x)=xy^k,\qquad \phi_k(y)=y
\]
is an automorphism for every \(k\in\mathbb Z\), by constructing its inverse.
\end{exercise}

\begin{exercise}
Apply \(\phi(x)=xy,\phi(y)=y\) to the word \(YXyX\), showing the image of
each capital letter before reducing.
\end{exercise}

\begin{exercise}
Let \(\eta(x)=x^2,\eta(y)=y\).  Explain why the universal property proves
that \(\eta\) is a homomorphism but not that it is an automorphism.  Prove
that it is not surjective by using exponent sums.
\end{exercise}

\begin{exercise}
Prove Theorem~\ref{thm:quotient-iso} in the special case
\(P=\langle x,y\mid x,y\rangle\) and
\(\phi(x)=xy,\phi(y)=y\), without invoking the general proof.
\end{exercise}

\section{Transporting AC moves through a fixed automorphism}
\label{sec:transport}

This is the central mathematical point.  The correct statement is about
reachability, not necessarily adjacency.

\begin{theorem}[Transported elementary move]
\label{thm:one-move}
Let \(\phi\in\Aut(\F)\).  If
\[
  (r_1,r_2)\xrightarrow{\text{one elementary AC move}}(r_1',r_2'),
\]
then
\[
  (\phi(r_1),\phi(r_2))
  \xrightarrow{\AC^*}
  (\phi(r_1'),\phi(r_2')),
\]
where \(\AC^*\) denotes a finite sequence of source-convention elementary
AC moves.
\end{theorem}

\begin{proof}
We inspect all three elementary moves.

\textbf{AC1.}  If \(r_i'=r_i r_j\), then
\[
  \phi(r_i')=\phi(r_i r_j)=\phi(r_i)\phi(r_j).
\]
Thus one transformed AC1 move suffices.

\textbf{AC2.}  If \(r_i'=r_i^{-1}\), then
\[
  \phi(r_i')=\phi(r_i^{-1})=\phi(r_i)^{-1}.
\]
Thus one transformed AC2 move suffices.

\textbf{AC3.}  Suppose \(r_i'=g r_i g^{-1}\), where \(g\) is one generator
or inverse.  Write the freely reduced word
\[
  \phi(g)=a_1a_2\cdots a_m,
  \qquad a_s\in\{x,X,y,Y\}.
\]
Starting from \(R=\phi(r_i)\), conjugate successively by
\(a_m,a_{m-1},\ldots,a_1\).  After \(m\) source-convention AC3 moves, the
result is
\[
  a_1\cdots a_m R a_m^{-1}\cdots a_1^{-1}
  =\phi(g)\phi(r_i)\phi(g)^{-1}
  =\phi(r_i').
\]
The other relator is untouched throughout.
\end{proof}

\subsection{Quantitative bounds}

\begin{definition}
For a fixed automorphism \(\phi\), define its generator expansion factor
\[
  \Lambda(\phi)
  =\max\{\len{\phi(x)},\len{\phi(y)}\}.
\]
Since inversion preserves reduced length,
\(\len{\phi(X)}=\len{\phi(x)}\) and
\(\len{\phi(Y)}=\len{\phi(y)}\).
\end{definition}

\begin{proposition}[Per-move transport bound]
\label{prop:per-move-bound}
Under the source-paper convention, the constructive transported costs are:
\[
\begin{array}{c|c}
\text{original elementary move}&
\text{number of transformed elementary moves}\\
\hline
\AC1&1\\
\AC2&1\\
\AC3\text{ by }g&\len{\phi(g)}\le\Lambda(\phi).
\end{array}
\]
These are guaranteed upper bounds for the displayed construction; they are
not claims about shortest AC distance.
\end{proposition}

\begin{corollary}[Whole-path bound]
\label{cor:path-bound}
Suppose an AC path contains \(N_1\) AC1 moves, \(N_2\) AC2 moves, and
AC3 conjugations by letters \(g_1,\ldots,g_{N_3}\).  Transporting the path
through \(\phi\) gives a path of length at most
\[
  N_1+N_2+\sum_{t=1}^{N_3}\len{\phi(g_t)}
  \le N_1+N_2+\Lambda(\phi)N_3.
\]
If \(N=N_1+N_2+N_3\), then the coarser bound is
\[
  N_\phi\le\Lambda(\phi)N.
\]
\end{corollary}

The last inequality uses \(\Lambda(\phi)\ge1\), which holds because an
automorphism cannot send a free generator to the identity.

\begin{example}
For \(\phi(x)=xy,\phi(y)=y\), one has \(\Lambda(\phi)=2\).
Conjugation by \(x\) or \(X\) costs at most two transformed AC3 moves;
conjugation by \(y\) or \(Y\) costs one.  Hence a path of length \(N\)
transports to one of length at most \(2N\), before correcting its terminal
basis.
\end{example}

\begin{proposition}[No automorphism-independent constant]
\label{prop:no-uniform}
There is no constant \(C\) such that every elementary AC3 move transports
through every \(\phi\in\Aut(\F)\) using at most \(C\) source-convention AC3
moves via the construction above.
\end{proposition}

\begin{proof}
Use \(\phi_k(x)=xy^k,\phi_k(y)=y\).  It is an automorphism with inverse
\(x\mapsto xy^{-k},y\mapsto y\).  A conjugation by \(x\) transports to a
conjugation by \(xy^k\), whose freely reduced length is \(k+1\) for
\(k\ge0\).  The constructive source-convention expansion therefore has
\(k+1\) AC3 moves, which is unbounded as \(k\to\infty\).
\end{proof}

\begin{remark}
Proposition~\ref{prop:no-uniform} concerns a uniform bound over all
automorphisms.  For each explicitly given \(\phi\), the bound
\(\Lambda(\phi)\) is finite and directly computable.  This effectiveness is
fundamentally different from the unbounded normal-closure parameter in
Lemma 11 of the stable-AC discussion.
\end{remark}

\subsection{From the transformed endpoint back to \((x,y)\)}

If an original path ends at \((x,y)\), its transported path ends at
\[
  (\phi(x),\phi(y)),
\]
not necessarily at \((x,y)\).  Because \(\phi\) is an automorphism, this pair
is a free basis.

\begin{theorem}[Nielsen generation in rank two]
\label{thm:nielsen-generation}
The group \(\Aut(\F)\) is generated by the elementary basis operations
\[
(u,v)\mapsto(v,u),\qquad
(u,v)\mapsto(u^{-1},v),\qquad
(u,v)\mapsto(uv,v),
\]
together with the corresponding operations on the second entry.
\end{theorem}

This is the classical Nielsen generation theorem.  On a relator pair, its
three displayed generators are respectively relator-order bookkeeping,
one AC2, and one AC1.  Hence a decomposition of \(\phi^{-1}\) into Nielsen
generators gives an explicit source-convention AC sequence from
\((\phi(x),\phi(y))\) to \((x,y)\).  This is stronger than merely observing
that the terminal pair generates \(\F\).

Let \(C_{\mathrm{basis}}(\phi)\) denote the number of source-convention AC
steps obtained from one chosen, explicitly tracked decomposition of
\(\phi^{-1}\) as in Theorem~\ref{thm:nielsen-generation}; swap steps cost
zero.  Combining it with
Corollary~\ref{cor:path-bound} gives
\[
  N_{\mathrm{total}}
  \le
  N_1+N_2+\Lambda(\phi)N_3+C_{\mathrm{basis}}(\phi).
\]
No claim is made that this is minimal.

The current project certificates perform an untracked Boolean Nielsen basis
check: they prove that the endpoint is a basis but do not store a
decomposition or a numerical value of \(C_{\mathrm{basis}}(\phi)\).
Obtaining a numerical endpoint bound therefore requires separately logging
a Nielsen reduction.  If that reduction uses the common replacements below,
they expand constructively as follows:
\[
\begin{array}{c|cccc}
\text{replacement of }a&
a b&a b^{-1}&b^{-1}a&b a\\
\hline
\text{source-convention cost}&1&3&3&5.
\end{array}
\]
For example, \(b^{-1}a=(a^{-1}b)^{-1}\) costs AC2, AC1, AC2; the other
three bounds follow by the same inversion-and-restoration constructions.
Any terminal signs require at most two additional AC2 moves.

\begin{example}[Endpoint cost in the source convention]
For \((\phi(x),\phi(y))=(xy,y)\), the tempting operation
\(xy\mapsto xy\,y^{-1}=x\) is not one source-convention AC1 move.  It is
the three-move sequence
\[
(xy,y)\xrightarrow{\AC2}(xy,Y)
\xrightarrow{\AC1}(xyY,Y)=(x,Y)
\xrightarrow{\AC2}(x,y).
\]
Thus this chosen source-convention endpoint correction has cost \(3\).
Under a macro convention allowing multiplication by \(r_2^{-1}\), it would
be encoded as one macro-move.
\end{example}

\begin{theorem}[Automorphism invariance of AC-triviality]
\label{thm:main-invariance}
For every balanced two-generator presentation \(P\) and every
\(\phi\in\Aut(\F)\),
\[
  P\text{ is AC-trivial}
  \quad\Longleftrightarrow\quad
  \phi(P)\text{ is AC-trivial}.
\]
\end{theorem}

\begin{proof}
Assume \(P\) is AC-trivial.  Transport its finite AC path to obtain
\[
  \phi(P)\xrightarrow{\AC^*}(\phi(x),\phi(y)).
\]
Choose a finite Nielsen-generator decomposition of \(\phi^{-1}\); it gives
a finite AC sequence from the terminal basis to \((x,y)\).  Hence
\(\phi(P)\) is AC-trivial.  Conversely, apply the same argument to
\(\phi^{-1}\).
\end{proof}

\[
\begin{array}{ccc}
P&\xrightarrow{\AC^*}&(x,y)\\
\downarrow\phi&&\downarrow\phi\\
\phi(P)&\xrightarrow{\AC^*}&(\phi(x),\phi(y))
\xrightarrow{\AC^*}(x,y).
\end{array}
\]

\begin{warning}[What the theorem does not say]
It does not prove \(P\sim_{\AC}\phi(P)\) for arbitrary \(P\).  An attempted
chain
\[
P\sim_{\AC}(x,y)\sim_{\AC}(\phi(x),\phi(y))
\sim_{\AC}\phi(P)
\]
already assumes the unknown first link.  The theorem says that \(P\) and
\(\phi(P)\) have the same truth value for AC-triviality; it does not say
that they lie in the same ordinary AC component.
\end{warning}

\subsection*{Exercises}

\begin{exercise}
For \(\phi(x)=xy,\phi(y)=y\), transport the AC3 move
\[
  r_1\mapsto x r_1X
\]
step by step, using only generator conjugations.
\end{exercise}

\begin{exercise}
An original path has 10 AC1 moves, 4 AC2 moves, 7 conjugations by \(x\),
and 3 conjugations by \(y\).  For the automorphism
\(\phi(x)=xyY,\phi(y)=Yx\), compute the direct bound from
Corollary~\ref{cor:path-bound}, after freely reducing the generator images.
\end{exercise}

\begin{exercise}
Prove that one source-convention AC1 and one AC2 always transport to one
move of the same type, even if substantial free cancellation occurs in
the displayed words.
\end{exercise}

\begin{exercise}
Explain precisely why Theorem~\ref{thm:one-move} should use
\(\xrightarrow{\AC^*}\) rather than a single arrow under the source-paper
convention, but may use a single arrow for AC3 under an arbitrary-word
conjugation convention.
\end{exercise}

\begin{exercise}
Find the logical error in the statement:
``Since \(P\) and \(\phi(P)\) present isomorphic trivial groups, they are
AC-equivalent.''
\end{exercise}

\section{Canonicalization and the eight signed relabelings}
\label{sec:canonical-relabel}

\subsection{Canonicalization}

The search stores a canonical representative of a relator pair.  For each
relator it:
\begin{enumerate}
  \item freely and cyclically reduces it;
  \item considers the relator and its inverse;
  \item selects the lexicographically least cyclic rotation;
  \item sorts the two canonical relators by length and lexical order.
\end{enumerate}

These operations do not invent an equivalence.

\begin{itemize}
  \item Free reduction preserves the element of \(\F\).
  \item Removing an outer inverse pair \(a w a^{-1}\mapsto w\) is one
  source-convention AC3 move, followed by free reduction.
  \item Whole-relator inversion is AC2.
  \item If \(r=uv\), the rotation \(uv\mapsto vu\) is conjugation:
  \[
    vu=u^{-1}(uv)u.
  \]
  Under generator-only AC3 it may require several elementary conjugations.
  \item Relator order is bookkeeping.
\end{itemize}

Canonicalization is therefore safe for AC-triviality, although it should
not be assigned zero elementary AC cost when one is explicitly counting
source-convention moves.

\subsection{The eight signed permutations}

A signed permutation sends \(x\) to one of \(x,X,y,Y\) and \(y\) to a
signed version of the other underlying generator.  There are
\[
  2!\,2^2=8
\]
such maps:
\[
\begin{array}{llll}
(x,y),&(X,y),&(x,Y),&(X,Y),\\
(y,x),&(Y,x),&(y,X),&(Y,X).
\end{array}
\]
Each is an automorphism whose inverse is another signed permutation.

For every input pair \(P\), the relabeling key is
\[
  K_{\mathrm{signed}}(P)
  =\min_{\sigma\in\Sigma_8}\canon(\sigma(P)).
\]
Equal keys mean the two presentations differ by a signed change of free
basis plus canonical bookkeeping.  By
Theorem~\ref{thm:main-invariance}, they have the same AC-triviality status.

Signed permutations are particularly well behaved computationally:
they preserve word length and commute with inversion and cyclic rotation.
Consequently their child sets under the encoded S-move agree after
relabeling.  General automorphisms such as \(x\mapsto xy\) do not have this
property because they change word lengths and rotation positions.

\subsection{Worked example: \(19\_45\) and \(19\_50\)}

The presentations have canonical relator pairs
\begin{align*}
19\_45&=(\texttt{YYYxxyyX},\ \texttt{YYYYXyxxYxx}),\\
19\_50&=(\texttt{YYYxyyXX},\ \texttt{YYYYXXYXXyx}).
\end{align*}
Use the signed automorphism
\[
  \sigma(x)=x,\qquad \sigma(y)=Y.
\]
For the first relator:
\[
\texttt{YYYxxyyX}
\xmapsto{\sigma}
\texttt{yyyxxYYX}
\xmapsto{\text{invert}}
\texttt{xyyXXYYY}
\xmapsto{\text{right rotate }3}
\texttt{YYYxyyXX}.
\]
For the second:
\[
\texttt{YYYYXyxxYxx}
\xmapsto{\sigma}
\texttt{yyyyXYxxyxx}
\xmapsto{\text{invert}}
\texttt{XXYXXyxYYYY}
\xmapsto{\text{right rotate }4}
\texttt{YYYYXXYXXyx}.
\]
Thus
\[
\canon(\sigma(19\_45))=19\_50.
\]
No search is needed.  The substitution is the mathematical change of
variables; inversion and rotation explain why the canonical strings do not
match immediately.

\subsection{Counts}

The 261 exact input pairs remain 261 after canonicalization: there are no
root collisions.  Quotienting by the eight signed permutations yields
\[
  261\longrightarrow171.
\]
This identifies 90 redundant rows.

\subsection*{Exercises}

\begin{exercise}
List the eight signed permutations and write the inverse of each one.
\end{exercise}

\begin{exercise}
Repeat the first-relator computation for \(19\_45\to19\_50\) without
looking at the worked derivation.  Explain why simply changing the case of
\(y\) does not immediately produce the stored target.
\end{exercise}

\begin{exercise}
Prove that signed permutations commute with free reduction and cyclic
rotation.  Identify the step that fails for \(\phi(x)=xy,\phi(y)=y\).
\end{exercise}

\begin{exercise}
If a signed-relabeling class has eight members and one representative is
proved AC-trivial in 200 moves, what can be concluded about the other seven?
What cannot be concluded about their shortest path lengths?
\end{exercise}

\section{Full \(\Aut(F_2)\) and Whitehead canonicalization}
\label{sec:whitehead}

The eight signed permutations form only a small subgroup of
\(\Aut(\F)\).  Length-changing Nielsen transformations, such as
\[
  x\mapsto xy,\qquad y\mapsto y,
\]
can identify presentations missed by signed relabeling.

\subsection{The twenty elementary Whitehead generators}

The implementation does not invent thousands of ad hoc maps.
\sourcefile{experiments/equivalence_classes/lib/autcanon.py}
always tries the same finite list of \textbf{20} elementary Whitehead
automorphisms.  Peak reduction and the length-preserving level-set search
\emph{compose} these maps; a certificate may therefore record a composite
such as \(x\mapsto\texttt{xxxxxxxy}\), \(y\mapsto X\), which is not itself one
of the twenty generators.

\paragraph{Eight first-kind maps (signed permutations).}
\[
\begin{array}{llll}
(x,y),&(x,Y),&(X,y),&(X,Y),\\
(y,x),&(y,X),&(Y,x),&(Y,X).
\end{array}
\]

\paragraph{Twelve second-kind maps (length-changing).}
\[
\begin{array}{lll}
x\mapsto x,\ y\mapsto yx &
x\mapsto x,\ y\mapsto Xy &
x\mapsto x,\ y\mapsto Xyx \\
x\mapsto x,\ y\mapsto yX &
x\mapsto x,\ y\mapsto xy &
x\mapsto x,\ y\mapsto xyX \\
x\mapsto xy,\ y\mapsto y &
x\mapsto Yx,\ y\mapsto y &
x\mapsto Yxy,\ y\mapsto y \\
x\mapsto xY,\ y\mapsto y &
x\mapsto yx,\ y\mapsto y &
x\mapsto yxY,\ y\mapsto y.
\end{array}
\]
Only the twelve second-kind maps are used in the length-descent loop; the eight
signed permutations already preserve total length.

\subsection{The classical theorem used}

\begin{theorem}[Whitehead peak reduction, specialized statement]
\label{thm:whitehead}
For a finite tuple of cyclic words in a finitely generated free group:
\begin{enumerate}
  \item if the tuple is not of minimum total length in its automorphism
  orbit, some Whitehead automorphism strictly decreases its total length;
  \item all minimum-length representatives in the same orbit are connected
  through length-preserving Whitehead moves.
\end{enumerate}
\end{theorem}

This classical theorem is cited rather than reproved.  Its hypotheses apply
to the object used here: an unordered pair of cyclic relators, with each
relator considered up to cyclic rotation and inversion.  The function
\(\canon\) chooses a word-pair representative of this object.

The project converts Theorem~\ref{thm:whitehead} into a complete orbit key:
\begin{enumerate}
  \item repeatedly apply a length-decreasing Whitehead automorphism;
  \item once minimum total length is reached, explore the entire
  length-preserving minimum level set;
  \item choose the lexicographically least canonical pair in that set.
\end{enumerate}

\begin{proposition}
For canonical relator pairs \(P,Q\), their Whitehead keys satisfy
\[
K(P)=K(Q)
\quad\Longleftrightarrow\quad
\exists\phi\in\Aut(\F)\text{ such that }
\canon(\phi(P))=\canon(Q).
\]
\end{proposition}

\begin{proof}
An orbit of unordered cyclic-relator pairs modulo inversion has one
minimum-level set.  By
Theorem~\ref{thm:whitehead}, every descent reaches its minimum length and
every minimum representative lies in the connected level set explored in
step 2.  Hence the lexicographic minimum is independent of the starting
canonical pair.  Equal keys therefore imply the displayed canonicalized
automorphism relation, while that relation necessarily produces the same
key.
\end{proof}

Peak reduction alone is not a canonical key: distinct descent choices can
reach different minimum representatives.  On the 261 inputs, using only
the peak-reduced endpoint would incorrectly produce 259 classes rather
than 168.  Exploring the complete minimum level set is essential.

\subsection{Independent checks}

The exact partition was computed independently by:
\begin{itemize}
  \item \sourcefile{experiments/analysis/whitehead.py}, and
  \item \sourcefile{experiments/equivalence_classes/lib/autcanon.py}.
\end{itemize}
They agree element for element.  The second implementation records a
witness \(\phi\) satisfying
\[
  \canon(\phi(P))=\text{reported orbit representative}.
\]
The certificate verifier then proves independently, by Nielsen reduction,
that every recorded pair \((\phi(x),\phi(y))\) is a basis.  Thus a bug that
produced a non-invertible substitution would be rejected.

\subsection{Empirical Aut witnesses on the 93 pure change-of-variable edges}

The shipped proof book
\sourcefile{results/equivalence_classes/certificates.json}
contains \(137\) verified edges in total:
\[
  93\ \texttt{cv}/\texttt{aut}
  \quad+\quad
  44\ \texttt{ac}/\texttt{aca}
  \quad=\quad
  137.
\]
The \(93\) edges are \textbf{pure automorphisms}: a single substitution \(\psi\)
with \(\canon(\psi(A))=\canon(B)\) and \emph{no} Andrews--Curtis move.  The
\(44\) remaining edges are the ones that use Definition~2.1 S-moves (ACA
search); they are not counted in the table below.

Among those \(93\) pure Aut edges, the witnessing substitution \(\psi\) has
the following empirical frequencies:

\begin{center}
\begin{tabular}{r l l}
\toprule
Count & Map \(\psi\) & Notes\\
\midrule
82 & \(x\mapsto x,\ y\mapsto Y\) & dominates; signed\\
3 & \(x\mapsto Y,\ y\mapsto x\) & signed\\
2 & \(x\mapsto y,\ y\mapsto x\) & signed\\
1 & \(x\mapsto X,\ y\mapsto y\) & signed\\
1 & \(x\mapsto xy,\ y\mapsto y\) & second-kind / composite\\
1 & \(x\mapsto xY,\ y\mapsto y\) & second-kind / composite\\
1 & \(x\mapsto XY,\ y\mapsto y\) & longer word\\
1 & \(x\mapsto xY,\ y\mapsto Y\) & mixed\\
1 & \(x\mapsto yx,\ y\mapsto Y\) & mixed\\
\bottomrule
\end{tabular}
\end{center}

Thus \(88\) of the \(93\) pure Aut merges are single-letter (signed-permutation
style), and only \(5\) need a longer generator image.  The textbook example
\(x\mapsto xy,\ y\mapsto y\) appears, but only once as a direct \(\texttt{cv}\)
witness; signed \(y\mapsto Y\) accounts for almost all Aut collapsing of the
\(261\) inputs.

\subsection{Worked example: \(19\_37\) and \(17\_41\)}

Take
\begin{align*}
19\_37&=(\texttt{YYXXyx},\ \texttt{YXyXyXyxxYxYx}),\\
17\_41&=(\texttt{YYXYXyx},\ \texttt{YXYxxxyXXX}).
\end{align*}
Apply \(\phi(x)=xy,\phi(y)=y\), proved invertible in
Example~\ref{ex:basic-aut}.  The proof book records:
\[
\texttt{YYXXyx}
\xmapsto{\phi}
\texttt{YYYXYXyxy}
\xmapsto{\text{reduce}}
\texttt{YYXYXyx},
\]
and
\[
\begin{aligned}
w_0&=\texttt{YXyXyXyxxYxYx},\\
w_1&=\texttt{YYXXXyxyxxxy}
  &&\text{(apply \(\phi\))},\\
w_2&=\texttt{YXXXyxyxxx}
  &&\text{(freely reduce)},\\
w_3&=\texttt{XXXYXYxxxy}
  &&\text{(invert)},\\
w_4&=\texttt{YXYxxxyXXX}
  &&\text{(right rotate by \(7\))}.
\end{aligned}
\]
Therefore
\[
  \canon(\phi(19\_37))=17\_41.
\]
This is a full-\(\Aut(\F)\) identification not restricted to one-letter
signed relabeling.

\subsection{Exact count}

Whitehead canonicalization gives
\[
  171\longrightarrow168.
\]
Unlike the later search-dependent numbers, 168 is exact with respect to
change of variables: no further \(\Aut(\F)\) orbit mergers exist among the
261 inputs.

\subsection*{Exercises}

\begin{exercise}
Explain why the map \(x\mapsto xy,y\mapsto y\) can change which numerical
rotations exist even though it preserves AC-triviality.
\end{exercise}

\begin{exercise}
Why does a witness satisfying
\(\canon(\phi(P))=\canon(Q)\) fail to prove anything unless \(\phi\) is
verified to be an automorphism?
\end{exercise}

\begin{exercise}
Give a concise proof that equality of complete Whitehead canonical keys is
an equivalence relation.  Which part depends on
Theorem~\ref{thm:whitehead}?
\end{exercise}

\begin{exercise}
Explain why agreement of two independently written implementations is
stronger evidence than running one implementation twice, but is not itself
a replacement for the mathematical theorem.
\end{exercise}

\section{Definition 2.1 S-moves as elementary AC sequences}
\label{sec:smove}

The equivalence search does not expand elementary AC1--AC3 one at a time.
It uses the solver's substitution or S-move encoding:
\[
  r_i\longmapsto
  \rot_{k_1}(r_i)\,
  \rot_{k_2}(r_j^\eps),
  \qquad i\ne j,\quad \eps\in\{+1,-1\},
\]
followed by reduction and canonicalization.  The other relator must remain
unchanged in the mathematical presentation.

\subsection{Rotations are conjugates}

Let \(r=uv\), and let the rotation move the suffix \(v\) to the front:
\[
  \rot(r)=vu.
\]
There are two useful conjugation identities:
\[
  vu=u^{-1}(uv)u=v(uv)v^{-1}.
\]
Thus one may choose either conjugator \(u^{-1}\) or \(v\).  Under
generator-only AC3, a conjugator \(c=a_1\cdots a_m\) expands into \(m\)
generator conjugations.  Define \(C(r,k)\) to be the length of the chosen
freely reduced conjugator realizing \(\rot_k(r)\).  A convenient bound is
\[
  C(r,k)\le\min\{k,\len{r}-k\}
\]
for a right rotation by \(k\), before any further cancellation.

\subsection{A constructive elementary realization}

Let
\[
R_i=c_1r_ic_1^{-1}=\rot_{k_1}(r_i),\qquad
R_j=c_2r_j^\eps c_2^{-1}=\rot_{k_2}(r_j^\eps),
\]
and put \(m_1=\len{c_1}\), \(m_2=\len{c_2}\).
The following source-convention sequence replaces \(r_i\) by \(R_iR_j\)
and restores the other relator to \(r_j\):

\begin{enumerate}
  \item If \(\eps=-1\), invert \(r_j\) using AC2.
  \item Conjugate \(r_i\) by \(c_1\), using \(m_1\) generator AC3 moves.
  \item Conjugate the current \(r_j^\eps\) by \(c_2\), using \(m_2\)
  generator AC3 moves.
  \item Apply one AC1 to replace the current target by its product with
  the current other relator.
  \item Conjugate the other relator by \(c_2^{-1}\), using \(m_2\) AC3
  moves, restoring it to \(r_j^\eps\).
  \item If \(\eps=-1\), invert the other relator again, restoring \(r_j\).
\end{enumerate}

\begin{proposition}[Source-convention S-move bound]
\label{prop:smove-bound}
Before final canonicalization, the construction above uses at most
\[
  \boxed{
  1+m_1+2m_2+2\,\mathbf 1_{\{\eps=-1\}}
  }
\]
elementary source-convention moves.
\end{proposition}

\begin{proof}
There is one AC1, \(m_1+2m_2\) generator AC3 moves, and two AC2 moves
exactly when \(\eps=-1\).  Inspection of the six stages proves that the
target is \(R_iR_j\) and the other relator is restored.
\end{proof}

This is a constructive upper bound, not a shortest-path formula.  Other
realizations may be shorter.  If arbitrary-word conjugation is counted as
one AC3 macro, the corresponding construction uses at most one AC1, three
arbitrary-word conjugations (target, other, restore other), and two
inversions in the negative case.  This explains why different papers quote
different expansion counts.

\subsection{Canonicalization cost}

The encoded child is then freely/cyclically reduced, possibly inverted,
rotated, and swapped.  To convert the raw child to that exact stored pair
under the source convention, add:
\begin{itemize}
  \item zero for free reduction and relator sorting;
  \item one AC2 for each whole-relator inversion selected;
  \item one generator AC3 per letter in each chosen cyclic-reduction or
  rotation conjugator.
\end{itemize}
Therefore the exact certificate supplies enough information to compute a
finite elementary upper bound, but the encoded S-move count itself is not
the elementary AC count.

\subsection{Transporting an S-move through \(\phi\)}

One may first expand the S-move into elementary operations and then apply
Proposition~\ref{prop:per-move-bound}.  Equivalently, replace each
conjugator \(c\) by \(\phi(c)\).  Since
\[
  \len{\phi(c)}\le\Lambda(\phi)\len{c},
\]
the pre-canonicalization transported construction has the bound
\[
  1+\len{\phi(c_1)}+2\len{\phi(c_2)}
  +2\,\mathbf 1_{\{\eps=-1\}}
\]
and hence
\[
  1+\Lambda(\phi)(m_1+2m_2)
  +2\,\mathbf 1_{\{\eps=-1\}}.
\]
Free cancellation can only lower these lengths.

\begin{warning}[Why child sets do not commute with general automorphisms]
The raw algebraic identity transports correctly, but the solver enumerates
rotations by numerical positions in the current words.  A length-changing
automorphism changes those positions and may introduce cancellation.
Therefore, for general \(\phi\),
\[
  \operatorname{children}(\phi(P))
  \ne
  \phi(\operatorname{children}(P))
\]
as enumerated canonical sets.  This does not threaten soundness: every
edge found is still legal.  It causes incompleteness because some quotient
edges may be missed when only one representative of an Aut-orbit is
expanded.
\end{warning}

\subsection*{Exercises}

\begin{exercise}
For \(r=abcdef\), express the right rotation by \(2\), \(efabcd\), as
conjugation in two different ways.  Compare the two conjugator lengths.
\end{exercise}

\begin{exercise}
Use Proposition~\ref{prop:smove-bound} when
\(m_1=2,m_2=3,\eps=-1\).  Give the number of AC1, AC2, and AC3 moves
separately.
\end{exercise}

\begin{exercise}
Write all intermediate relator pairs in the six-stage construction when
\(c_1=x\), \(c_2=Y\), and \(\eps=-1\).  Do not simplify away the restoration
steps.
\end{exercise}

\begin{exercise}
Let \(\phi(x)=xy,\phi(y)=y\).  If \(c_1=Xy\) and \(c_2=xY\), calculate
\(\phi(c_1),\phi(c_2)\), freely reduce them, and evaluate the sharper
transported S-move bound for \(\eps=+1\).
\end{exercise}

\section{ACA: the equivalence relation used by the search}
\label{sec:aca}

\begin{definition}
Write \(P\sim_{\ACA}Q\) when \(P\) and \(Q\) are connected by a finite
sequence whose steps are either:
\begin{enumerate}
  \item ordinary AC moves, or
  \item changes of variables by automorphisms in \(\Aut(\F)\).
\end{enumerate}
\end{definition}

An automorphism step is not declared to be an AC step.  Instead, the
important Boolean property
\[
  \mathcal T(P)=
  \begin{cases}
  1,&P\text{ is AC-trivial},\\
  0,&P\text{ is not AC-trivial}
  \end{cases}
\]
is invariant under each type of ACA step.

\begin{proposition}[ACA preserves the problem]
\label{prop:aca-preserves}
If \(P\sim_{\ACA}Q\), then
\[
  P\text{ is AC-trivial}
  \quad\Longleftrightarrow\quad
  Q\text{ is AC-trivial}.
\]
\end{proposition}

\begin{proof}
An AC move preserves membership in the AC component of \((x,y)\), because
it is reversible by a finite AC sequence.  An automorphism preserves
AC-triviality by Theorem~\ref{thm:main-invariance}.  Equality of the truth
value propagates along a finite chain.
\end{proof}

\begin{warning}
``Same ACA problem'' is weaker than ``joined by a raw AC path.''  The search
therefore reports components of a finite verified ACA-edge subgraph, not
proven ordinary AC-equivalence classes and not an exact partition into full
ACA classes.
\end{warning}

\subsection{The quotient search graph}

The search uses:
\begin{itemize}
  \item \textbf{Node:} one complete \(\Aut(\F)\)-orbit, keyed by its
  Whitehead canonical representative.
  \item \textbf{Edge:} apply one encoded S-move to the chosen Aut-minimal
  representative, then map the child to its Whitehead key.
  \item \textbf{Priority:} the Aut-minimal total relator length.
\end{itemize}

Using Aut-minimal length matters.  A raw word pair of total length 34 may be
one change of variables away from length 16.  A raw length-ordered search
must enumerate a large ``hump'' that the quotient representation collapses.

\subsection{Soundness and incompleteness}

\begin{theorem}[Soundness of every reported merge]
\label{thm:search-sound}
If two sources meet in the quotient search, they have the same
AC-triviality status.
\end{theorem}

\begin{proof}
Each path step consists of a legal S-move, which
Proposition~\ref{prop:smove-bound} expands into elementary AC moves,
followed by a verified automorphism and canonical bookkeeping.  Each part
preserves \(\mathcal T\).  If paths from \(P\) and \(Q\) end in the same
Aut-orbit, transitivity gives \(\mathcal T(P)=\mathcal T(Q)\).
\end{proof}

The search is incomplete for three reasons:
\begin{enumerate}
  \item it expands only one representative from each Aut-orbit;
\item general automorphisms do not preserve the enumerated child set;
\item length, state, time, and memory caps remove possible states.
\end{enumerate}
All three mechanisms remove edges; none can create a false edge.  Hence:
\[
\boxed{\text{Every merge found is a proof.  No merge found proves nothing.}}
\]

\subsection*{Exercises}

\begin{exercise}
Construct a three-step ACA chain containing one automorphism step and two
AC steps.  Explain why it proves equality of AC-triviality status without
proving that the endpoints are directly adjacent in the AC graph.
\end{exercise}

\begin{exercise}
Why does expanding only one representative of an Aut-orbit preserve
soundness but lose completeness?
\end{exercise}

\begin{exercise}
Suppose a search cap removes the only path connecting two reported
components.  Does this create too many components or too few?  Why is this
the safe direction for choosing a workload?
\end{exercise}

\begin{exercise}
State the precise difference between a node of the raw AC graph and a node
of the ACA quotient graph.
\end{exercise}

\section{The complete reduction: 261 to at most 124}
\label{sec:pipeline}

\begin{center}
\begin{tabularx}{\textwidth}{X r X}
\toprule
Identification applied & Count & Status\\
\midrule
Exact ordered relator pair & 261 & exact\\
Canonical rotation/inversion/swap & 261 & exact; no collisions\\
Eight signed generator relabelings & 171 & exact for this subgroup\\
Full \(\Aut(\F)\), Whitehead canonical key & 168 & exact Aut-orbit partition\\
ACA search, Aut-minimal ceiling 28 & 126 & sound upper bound\\
Probe at ceiling 34 & 125 & sound upper bound\\
Overnight ladder, ceilings 30--36 & 124 & sound upper bound\\
\bottomrule
\end{tabularx}
\end{center}

\subsection{From 168 to 126}

The canonical cap-28 sweep seeded all 261 roots and the trivial
presentation.  Initial Aut coincidences supplied 93 change-of-variable
forest edges, reducing 261 to 168.  The ACA search supplied 42 further
verified edges, reducing
\[
  168\longrightarrow126.
\]
The resulting 135 edges form a spanning forest on the 261 input vertices:
\[
  261-135=126.
\]
This historical 126-component forest is recorded in
\sourcefile{results/equivalence_classes/sweep_seam_28_250.json}.  The
current proof book and certificates were regenerated after adding the two
later edges and contain 124 components.

\subsection{Why 126 became 125}

The cap-28 search could not reach an Aut-minimal meeting state of total
length 30, regardless of node budget.  A probe with ceiling 34 found
\[
  21\_3\sim_{\ACA}21\_29,
\]
joining two singleton components and giving 125.  The paths had lengths
one and three encoded S-moves to a common Aut-orbit.

\subsection{Why 125 became 124}

Five overnight arms at ceilings 30, 32, 32 (deterministic rerun), 34, and
36 all found the additional merge
\[
  21\_7\sim_{\ACA}21\_28.
\]
Again the two sides reached a common Aut-orbit in one and three encoded
S-moves.  This joins another pair of singleton components:
\[
  126-2=124.
\]
The five arms explored roughly 700,000 nodes in total and found no further
merger, but every arm ended because of a state or memory resource cap, not
because its queue was exhausted.

\subsection{The coverage theorem}

\begin{theorem}[One representative per certified component suffices]
\label{thm:coverage}
Let \(\mathcal P\) be the 261 input presentations, and let the verified ACA
edges have 124 connected components.  Choose one representative \(R_C\)
from each component \(C\).  If every \(R_C\) is AC-trivial, then every
presentation in \(\mathcal P\) is AC-trivial.
\end{theorem}

\begin{proof}
For every \(P\in\mathcal P\), there is exactly one certified component
\(C\) containing it and a finite verified edge path from \(P\) to \(R_C\).
By Proposition~\ref{prop:aca-preserves},
\[
  P\text{ is AC-trivial}
  \Longleftrightarrow
  R_C\text{ is AC-trivial}.
\]
The assumed triviality of all 124 representatives therefore implies the
triviality of all 261 inputs.
\end{proof}

\begin{corollary}[Safe-direction error]
If the incomplete search missed additional valid mergers, running 124
representatives performs redundant work but does not omit any of the 261
inputs.
\end{corollary}

\begin{proof}
A missed edge can only merge two currently separate certified components,
so the true ACA partition can be coarser and have fewer components.  It
cannot split a verified component.  Every original vertex is already in
one selected component.
\end{proof}

\subsection{How to obtain a 124-representative workload}

The historical probe-seed manifest
\sourcefile{results/equivalence_classes/classes_126_from_greedy_1000000_261_mrl48.jsonl}
still contains
126 rows and is intentionally retained because the probe certificates name
its class IDs.  Starting from that historical manifest, merge:
\[
\begin{array}{ccl}
C123\ (21\_29)&\text{with}&C124\ (21\_3),\\
C122\ (21\_28)&\text{with}&C125\ (21\_7).
\end{array}
\]
Keep either representative from each pair.  All other canonical
components remain unchanged, producing 124 representatives.  The current
production deliverable has already performed these unions:
\begin{center}
\sourcefile{results/equivalence_classes/classes_sweep_seam_28_250_plus_overnight.csv}.
\end{center}
It contains 124 rows and supersedes the historical 126-row workload.

\subsection{What 124 does and does not mean}

\begin{itemize}
  \item It means there are \emph{at most} 124 distinct ACA problems among
  these 261 inputs.
  \item It does not mean that the 124 certified components are mutually
  ACA-inequivalent; more search may merge them.
  \item It does not mean there are 124 ordinary AC-classes.
  \item It does not prove any input AC-trivial.  The reported solved set is
  empty.
  \item It proves a conditional workload statement: 124 successful
  trivializations would settle all 261.
\end{itemize}

\subsection*{Exercises}

\begin{exercise}
Verify the component arithmetic
\[
261-93-42-2=124.
\]
Explain what each subtraction represents mathematically.
\end{exercise}

\begin{exercise}
Suppose a future verified run finds five additional edges, but two lie
inside already connected components.  By how much can the component count
decrease?
\end{exercise}

\begin{exercise}
Give a formal counterexample to the statement:
``Because the search returned 124 components, it proved there are at least
124 distinct problems.''
\end{exercise}

\begin{exercise}
If only 20 of the selected 124 representatives are trivialized, what
exactly has been proved about the 261 inputs?
\end{exercise}

\begin{exercise}
Explain why five runs agreeing on 124 strengthens empirical confidence but
does not turn the upper bound into an exact count.
\end{exercise}

\section{Certificates: why the mergers are proofs}
\label{sec:certificates}

The search output is not trusted merely because it prints a class count.
The verification layers are deliberately separated.

\subsection{Two edge types}

The current proof book has exactly \(137\) edges.  They split as follows.

\begin{description}
  \item[\texttt{cv}/\texttt{aut} --- \(93\) edges.]
  \textbf{Pure Aut / change of variables only.}  A recorded automorphism
  \(\psi\) must satisfy
  \[
    \canon(\psi(A))=\canon(B).
  \]
  No Definition~2.1 S-move appears.  These \(93\) forest edges are exactly
  what collapses \(261\to168\) Aut-orbits.  Empirically, \(82\) of them use
  \(\psi\colon x\mapsto x,\ y\mapsto Y\); see the frequency table in
  Section~\ref{sec:whitehead}.

  \item[\texttt{ac}/\texttt{aca} --- \(44\) edges.]
  \textbf{ACA edges: AC/S-moves are required.}  Both roots have replayable
  Definition~2.1 paths, optionally interleaved with verified automorphisms,
  ending in the same Aut-orbit.  The historical cap-28 sweep found \(42\)
  such edges (\(168\to126\)); the two later probe mergers
  (\(21\_3\sim 21\_29\) and \(21\_7\sim 21\_28\)) bring the current book to
  \(44\).  Of these \(44\), only \(6\) are tagged pure AC paths (every
  recorded Aut witness on the path is the identity after Aut-minimal
  alignment); the rest mix nontrivial \(\phi\) with S-moves.
\end{description}

\[
  \underbrace{93}_{\text{Aut only}}
  \;+\;
  \underbrace{44}_{\text{ACA with S-moves}}
  \;=\;
  137
  \qquad\text{and}\qquad
  261-137=124.
\]

An \texttt{aca} edge generally does not satisfy
\(\canon(\psi(A))=\canon(B)\) for one direct substitution; AC moves occur
between the endpoints.  Applying the pure-change test to an \texttt{aca}
edge would reject a correct certificate.

\subsection{Independent verification}

The verifier performs the following checks.

\begin{enumerate}
  \item Re-read each source presentation from the original CSV, binding
  both its row identity and its relator words.
  \item Replay each Definition 2.1 move by direct word operations.
  \item Check every recorded \(\phi\) by independent Nielsen basis
  reduction.
  \item Canonicalize independently of the pure-Python search routine.
  \item Require both sides of an ACA merge to end in the same recorded
  Aut-orbit.
  \item Check the exponent-sum determinant invariant.
  \item Rebuild the entire partition using only edges that passed all
  previous checks.
\end{enumerate}

The last step is essential.  Checking each listed edge cannot detect a
separate union-find bug that assigns an unconnected vertex to a component.
Rebuilding the transitive closure from verified edges detects such an
over-merge.

The production verifier
\begin{center}
\sourcefile{experiments/equivalence_classes/verify/verify_proofs.py}.
\end{center}
re-proves the current 137-edge, 124-component proof book from
\sourcefile{results/equivalence_classes/certificates.json} and the raw input
CSV.  The two mergers that
historically changed 126 to 124 are also checked directly in their original
probe files by
\begin{center}
\sourcefile{experiments/equivalence_classes/verify/verify_probe_merges.py}.
\end{center}
The latter reconstructs roots from the historical 126 manifest rather than
trusting the probe files' own root blocks.

\subsection{Mutation testing}

The checker is tested by deliberately corrupting certificates: replacing
an automorphism with a non-basis endomorphism, changing a rotation number,
swapping words between presentation IDs, corrupting a concatenation, or
forcing an unsupported over-merge.  Each mutation must cause verification
to fail.  This tests not only the mathematical certificate but also whether
the checker actually enforces the advertised conditions.

\subsection{Invariant cross-check}

For a pair \(P=(r_1,r_2)\), form the exponent-sum matrix
\[
  M(P)=
  \begin{pmatrix}
  \exp_x(r_1)&\exp_y(r_1)\\
  \exp_x(r_2)&\exp_y(r_2)
  \end{pmatrix}.
\]
Elementary AC1 and AC2 perform unimodular row operations; AC3 does not
change exponent sums.  A change of free basis performs a unimodular column
operation.  Therefore
\[
  \len{\det M(P)}
\]
is ACA-invariant.  It equals 1 for every input, so it cannot separate the
remaining groups, but a mismatch across a claimed merge would expose an
error.

\subsection*{Exercises}

\begin{exercise}
Why is checking every edge insufficient to verify the reported component
count?
\end{exercise}

\begin{exercise}
Show directly how AC1, AC2, and AC3 act on the exponent-sum matrix and prove
that \(\len{\det M}\) is unchanged.
\end{exercise}

\begin{exercise}
Explain why the fact that \(\len{\det M}=1\) for all 261 is useful as a
bug detector but useless as a discriminator among these inputs.
\end{exercise}

\begin{exercise}
Design one certificate mutation not listed above and state the verifier
condition that should reject it.
\end{exercise}

\section{Presentation-ready synthesis}
\label{sec:presentation}

\subsection{A concise theorem chain}

\begin{enumerate}
  \item A balanced presentation
  \(\langle x,y\mid r_1,r_2\rangle\) is the quotient
  \(\F/\Ncl{r_1,r_2}\).
  \item An automorphism \(\phi\in\Aut(\F)\) is a reversible change of free
  basis and induces an isomorphism of the presented quotient groups.
  \item Every elementary AC move transports through \(\phi\) to a finite
  source-convention AC sequence.  AC1 and AC2 cost one; AC3 by \(g\) costs
  at most \(\len{\phi(g)}\).
  \item Hence \(P\) is AC-trivial if and only if \(\phi(P)\) is
  AC-trivial.
  \item Therefore each AC or Aut edge preserves the question ``is this
  presentation AC-trivial?''
  \item The finite verified ACA-edge subgraph on 261 inputs has exactly
  124 components, each contained in a true ACA class.
  \item Consequently one representative per component is a sufficient
  coverage set; proving all 124 representatives AC-trivial would prove all
  261 AC-trivial.
\end{enumerate}

\subsection{Move-bound summary}

\begin{center}
\begin{tabularx}{\textwidth}{X X}
\toprule
Object & Constructive source-convention bound\\
\midrule
Transported AC1 & \(1\)\\
Transported AC2 & \(1\)\\
Transported AC3 by \(g\) & \(\len{\phi(g)}\le\Lambda(\phi)\)\\
Transported path & \(N_1+N_2+\Lambda(\phi)N_3\), before basis correction\\
Terminal basis correction &
\(C_{\mathrm{basis}}(\phi)\), from a separately logged Nielsen reduction;
not stored by the current certificates\\
One raw S-move before canonicalization &
\(1+m_1+2m_2+2\mathbf1_{\{\eps=-1\}}\)\\
Transported raw S-move &
\(1+\Lambda(\phi)(m_1+2m_2)+2\mathbf1_{\{\eps=-1\}}\)\\
\bottomrule
\end{tabularx}
\end{center}

\subsection{Claims to make}

\begin{itemize}
  \item ``The 261 inputs represent at most 124 distinct ACA problems.''
  \item ``The 168 full-\(\Aut(\F)\) count is exact.''
  \item ``The 124 count is the transitive closure of certified mergers and
  is still an upper bound.''
  \item ``One representative per certified component covers all 261.''
  \item ``Zero of the 261 were trivialized by this analysis.''
\end{itemize}

\subsection{Claims not to make}

\begin{itemize}
  \item ``There are exactly 124 AC-classes.''
  \item ``The 124 certified components are proven ACA-inequivalent.''
  \item ``\(P\) is always AC-equivalent to \(\phi(P)\).''
  \item ``An automorphism preserves the same numerical rotation or the
  same elementary move count.''
  \item ``The solver's \(r_j^{-1}\) multiplication is one elementary AC1
  move under the source convention.''
  \item ``The equivalence-class computation solved any presentation.''
\end{itemize}

\subsection{Oral rehearsal questions}

\begin{exercise}
Give a two-minute explanation of why \(\phi(x)=xy,\phi(y)=y\) is an
automorphism and why it preserves AC-triviality.
\end{exercise}

\begin{exercise}
Give a one-minute explanation of why 124 is sufficient but not exact.
\end{exercise}

\begin{exercise}
An audience member asks: ``If ACA-equivalence is not AC-equivalence, why
are you allowed to discard 137 runs?''  Give a rigorous response.
\end{exercise}

\begin{exercise}
An audience member asks: ``Did you prove 124 presentations trivial?''
Answer without weakening or overstating the result.
\end{exercise}

\newpage
\section{Complete answer key and learning diagnostics}
\label{sec:answers}

The answers deliberately show intermediate reasoning.  When an exercise can
be solved in several ways, the answer gives one auditable construction rather
than claiming minimality.

\subsection*{Answers for Section 1: AC conventions}

\subsubsection*{Answer to Exercise 1.1}

Invert the second relator, multiply, and restore it:
\[
(r_1r_2,r_2)
\xrightarrow{\AC2}
(r_1r_2,r_2^{-1})
\xrightarrow{\AC1}
(r_1r_2r_2^{-1},r_2^{-1})
=(r_1,r_2^{-1})
\xrightarrow{\AC2}
(r_1,r_2).
\]
This also proves that the directed-looking positive AC1 operation is
reversible by a finite sequence, which is enough for AC-equivalence.

\diagnostic{If the middle step looked like multiplication by an inverse,
notice that the \emph{current} second relator is already \(r_2^{-1}\).  AC1
is still positive multiplication by the current relator.}

\subsubsection*{Answer to Exercise 1.2}

If AC3 permits \(r\mapsto w r w^{-1}\) for arbitrary \(w\), then any cyclic
rotation is one AC3 macro-move.  In the source-paper convention, AC3 permits
only one generator or inverse.  Conjugation by a reduced word
\(w=a_1\cdots a_m\) then expands into \(m\) elementary AC3 moves.  Thus
``one rotation'' may mean one encoded operation or several elementary
source-paper moves.

\diagnostic{Whenever a move count is quoted, first write the permitted
conjugator beside AC3.}

\subsubsection*{Answer to Exercise 1.3}

\begin{itemize}
  \item \(r_1\mapsto r_1^{-1}\) is one elementary AC2.
  \item \(r_1\mapsto r_1r_2^{-1}\), while restoring \(r_2\), is a
  three-move macro: AC2, AC1, AC2.
  \item \(r_1\mapsto w r_1w^{-1}\) is elementary only when the freely
  reduced \(w\) is one generator or inverse.  Otherwise it is an
  arbitrary-word macro that expands into \(\len{w}\) generator AC3 moves.
  \item \((r_1,r_2)\mapsto(r_2,r_1)\) is bookkeeping because relator order
  is not part of the presentation.
\end{itemize}

\diagnostic{The operation name alone is insufficient; check both the
allowed operand and what must remain unchanged.}

\subsection*{Answers for Section 2: free groups and presentations}

\subsubsection*{Answer to Exercise 2.1}

Reduce from left to right:
\[
\texttt{xyYXxXY}
\longrightarrow \texttt{xXxXY}
\longrightarrow \texttt{xXY}
\longrightarrow \texttt{Y}.
\]
More explicitly, \(yY\), then \(xX\), then the remaining \(xX\) cancel.
The reduced word is \(Y\), whose inverse is \(y\).

\diagnostic{If a different answer was obtained, use a stack: cancel only
adjacent inverse letters and re-check adjacency after each cancellation.}

\subsubsection*{Answer to Exercise 2.2}

The relation \(y=1\) and the relation \(xy=1\) imply \(x=1\), so the
presented group is trivial.  Under the source convention:
\[
(xy,y)\xrightarrow{\AC2}(xy,Y)
\xrightarrow{\AC1}(xyY,Y)=(x,Y)
\xrightarrow{\AC2}(x,y).
\]
Thus the presentation is not merely group-trivial; this explicit sequence
proves it AC-trivial.

\diagnostic{Separating the quotient argument from the AC sequence is
important: the first proves the group is trivial, while the second proves
the stronger AC-triviality statement.}

\subsubsection*{Answer to Exercise 2.3}

\(G_P=1\) says that \(r_1,r_2\) normally generate \(\F\).  A normal-closure
membership statement permits an unspecified product of conjugates and gives
no prescribed sequence reducing the \emph{ordered relator pair} through
AC1--AC3.  Asserting that such a path always exists is precisely the
Andrews--Curtis conjecture.

\diagnostic{Do not replace ``normally generates'' by ``has a known AC
derivation.''  They are different strengths of information.}

\subsubsection*{Answer to Exercise 2.4}

Balancedness means that the number of relators equals the number of
generators.  The standard AC conjecture compares balanced \(n\)-generator,
\(n\)-relator presentations to the \(n\)-generator trivial presentation.
The equivalence pipeline fixes \(n=2\) and stores two relator slots.  A
two-generator, three-relator presentation is not balanced and is neither a
vertex in that graph nor one of the 261 inputs.

\diagnostic{Stable AC may temporarily change the number of generators and
relators together, but stabilization is not used in the 261-to-124
pipeline.}

\subsection*{Answers for Section 3: automorphisms}

\subsubsection*{Answer to Exercise 3.1}

Define
\[
\psi_k(x)=xy^{-k},\qquad\psi_k(y)=y.
\]
Then
\[
\psi_k(\phi_k(x))=\psi_k(xy^k)=xy^{-k}y^k=x,
\qquad
\psi_k(\phi_k(y))=y.
\]
Similarly,
\[
\phi_k(\psi_k(x))=\phi_k(xy^{-k})=xy^ky^{-k}=x,
\qquad
\phi_k(\psi_k(y))=y.
\]
The two homomorphisms are inverse on the free generators and hence inverse
on all of \(\F\).

\diagnostic{It is enough to check compositions on \(x,y\) only because of
the universal property, not because longer words are being ignored.}

\subsubsection*{Answer to Exercise 3.2}

The forced letter images are
\[
\phi(Y)=Y,\qquad \phi(X)=YX,\qquad
\phi(y)=y,\qquad \phi(X)=YX.
\]
Therefore
\[
\phi(\texttt{YXyX})
=\texttt{Y}\,\texttt{YX}\,\texttt{y}\,\texttt{YX}
=\texttt{YYXyYX}.
\]
The adjacent \(yY\) cancels, leaving
\[
\boxed{\texttt{YYXX}}.
\]

\diagnostic{For a capital letter, invert the whole image word: reverse it
and invert every letter.  A simple case swap is wrong for multi-letter
images.}

\subsubsection*{Answer to Exercise 3.3}

The universal property gives an endomorphism because any two image words
define a homomorphism.  Let \(\exp_x(w)\) denote the exponent sum of \(x\).
For every word \(w\),
\[
\exp_x(\eta(w))=2\exp_x(w),
\]
which is even.  The element \(x\in\F\) has exponent sum \(1\), so it is not
in the image of \(\eta\).  Thus \(\eta\) is not surjective and cannot be an
automorphism.

\diagnostic{A determinant or exponent-sum test can disprove
invertibility, but passing an abelianized test is not generally enough to
prove a free-group endomorphism invertible.}

\subsubsection*{Answer to Exercise 3.4}

The transformed relators are \((xy,y)\).  Their normal closure is all of
\(\F\), because it contains \(y\) and
\[
(xy)y^{-1}=x.
\]
Hence the transformed quotient is trivial.  The explicit inverse
\(\psi(x)=xY,\psi(y)=y\) proves that the map induced by \(\phi\) on the
quotients has a two-sided inverse.  This is the special case of
Theorem~\ref{thm:quotient-iso}.

\diagnostic{The quotient proof uses normal closure; the later three-move
sequence proves the stronger AC statement.}

\subsection*{Answers for Section 4: transported moves and bounds}

\subsubsection*{Answer to Exercise 4.1}

The original target is \(xr_1X\).  Applying \(\phi\) gives
\[
\phi(xr_1X)=xy\,\phi(r_1)\,YX.
\]
Starting with \(R=\phi(r_1)\), first conjugate by \(y\), then by \(x\):
\[
R\xrightarrow{\AC3}yRY
\xrightarrow{\AC3}x(yRY)X=xyRYX.
\]
The second relator is unchanged.  Thus one original generator conjugation
became two source-convention AC3 moves.

\diagnostic{Conjugate by the letters of \(\phi(x)=xy\) in reverse order:
first \(y\), then \(x\).}

\subsubsection*{Answer to Exercise 4.2}

Freely reduce the images:
\[
\phi(x)=\texttt{xyY}\longrightarrow\texttt{x},
\qquad
\phi(y)=\texttt{Yx}.
\]
This is an automorphism; for example, its inverse is
\(\psi(x)=x,\psi(y)=xY\), since
\[
\psi(Yx)=yXx=y,\qquad \phi(xY)=xXy=y,
\]
and both compositions also fix \(x\).  The image lengths are \(1\) and \(2\), so
\(\Lambda(\phi)=2\).  The sharper sum is
\[
10+4+7\cdot1+3\cdot2
=14+7+6
=\boxed{27}.
\]
The coarser bound \(\Lambda(\phi)N\), with
\(N=10+4+7+3=24\), would be \(48\).

\diagnostic{Reduce \(\phi(x),\phi(y)\) before measuring them, and use the
actual conjugating letters for the sharper sum.}

\subsubsection*{Answer to Exercise 4.3}

For AC1,
\[
\phi(r_ir_j)=\phi(r_i)\phi(r_j)
\]
as a group identity.  Any cancellation merely changes the displayed
reduced representative after the one AC1 operation.  For AC2,
\[
\phi(r_i^{-1})=\phi(r_i)^{-1}.
\]
Again free reduction is equality in \(\F\), not an extra move.  Therefore
one transformed operation of the same type reaches the correct free-group
element.

\diagnostic{Move count concerns operations on group elements; free
reduction chooses their word representatives.}

\subsubsection*{Answer to Exercise 4.4}

Under generator-only AC3, \(\phi(g)\) may be a word of length greater than
one, so the conjugation must be decomposed letter by letter and the correct
arrow is \(\AC^*\).  If AC3 is defined to permit an arbitrary conjugating
word, the single operation
\[
\phi(r_i)\mapsto\phi(g)\phi(r_i)\phi(g)^{-1}
\]
is one AC3 macro-move.  AC1 and AC2 remain one move in either convention.

\diagnostic{This is a convention issue, not a failure of the homomorphism
identity.}

\subsubsection*{Answer to Exercise 4.5}

Isomorphic trivial quotient groups only show that both presentations
present the group \(1\).  The implication ``presents \(1\)'' \(\Rightarrow\)
``is AC-trivial'' is exactly the unproved Andrews--Curtis assertion.
The correct route is stronger: when one presentation already has an AC
path, transport that path and then reduce the terminal basis.

\diagnostic{Identify whether an argument proves a fact about the quotient
group or a path in the AC graph.}

\subsection*{Answers for Section 5: canonicalization and relabeling}

\subsubsection*{Answer to Exercise 5.1}

The four unswapped maps are self-inverse:
\[
(x,y),\quad(X,y),\quad(x,Y),\quad(X,Y).
\]
Among the swapped maps,
\[
(y,x)^{-1}=(y,x),\qquad
(Y,X)^{-1}=(Y,X),
\]
and
\[
(Y,x)^{-1}=(y,X),\qquad
(y,X)^{-1}=(Y,x).
\]
Here \((u,v)\) abbreviates \(x\mapsto u,y\mapsto v\).  Direct composition on
\(x,y\) verifies every claim.

\diagnostic{For a swapped signed map, solve for the preimages of \(x\) and
\(y\); do not assume every signed permutation is self-inverse.}

\subsubsection*{Answer to Exercise 5.2}

First substitute \(y\mapsto Y\), remembering \(Y\mapsto y\):
\[
\texttt{YYYxxyyX}\longmapsto\texttt{yyyxxYYX}.
\]
Invert by reversing and changing case:
\[
\texttt{yyyxxYYX}^{-1}=\texttt{xyyXXYYY}.
\]
Move the final three letters to the front:
\[
\texttt{xyyXXYYY}\longmapsto\texttt{YYYxyyXX}.
\]
This is the first relator of \(19\_50\).  The direct substituted string is
not stored because canonicalization may invert and rotate it.

\diagnostic{A canonical target is a fingerprint, not necessarily the raw
substitution output.}

\subsubsection*{Answer to Exercise 5.3}

A signed permutation maps each letter to one letter and maps inverse
letters to inverse letters.  Hence an adjacent inverse pair maps to an
adjacent inverse pair, so free reduction commutes.  Because lengths and
positions are unchanged, mapping a cyclic rotation equals rotating the
mapped word by the same position.

For \(\phi(x)=xy\), one \(x\)-position becomes two positions and cancellation
may occur at image boundaries.  The same numerical rotation need no longer
select the image of the same block.

\diagnostic{The decisive property is length preservation at the letter
level, not merely invertibility.}

\subsubsection*{Answer to Exercise 5.4}

All eight have the same AC-triviality truth value by
Theorem~\ref{thm:main-invariance}.  Therefore a 200-move proof for one
settles all seven others.  It does not imply that each has a shortest path
of length 200, or even that the mechanically transported path has length
200: conjugating letters and terminal basis corrections may change the
elementary count.

\diagnostic{Transfer the existence of a proof, not its optimal path
length.}

\subsection*{Answers for Section 6: full \(\Aut(F_2)\)}

\subsubsection*{Answer to Exercise 6.1}

An encoded rotation index refers to positions in the current word.  The
map \(x\mapsto xy\) expands every \(x\) to two letters and may create free
cancellation next to neighboring images.  Consequently word length and
the boundaries of cyclic blocks change.  Nevertheless, the underlying
rotation is a conjugation, and its image is conjugation by the image word,
so AC-triviality is preserved.

\diagnostic{Numerical rotation indices are representation-dependent;
conjugation identities are group-theoretic.}

\subsubsection*{Answer to Exercise 6.2}

An arbitrary endomorphism still satisfies the homomorphism identities, but
it may send the terminal basis \((x,y)\) to a non-basis.  Without
invertibility there is no guaranteed path from
\((\phi(x),\phi(y))\) back to \((x,y)\), nor a reverse implication obtained
from \(\phi^{-1}\).  Thus the entire AC-triviality equivalence can fail.

\diagnostic{The certificate must establish ``basis,'' not merely
``homomorphism.''}

\subsubsection*{Answer to Exercise 6.3}

Equality of canonical keys is reflexive, symmetric, and transitive because
ordinary equality is.  The substantive statement is that equality of keys
is equivalent to equality of Aut-orbits of unordered cyclic-relator pairs,
with each relator taken modulo inversion.  Whitehead peak reduction supplies
minimum length, and minimum-level connectivity proves that every member of
one orbit chooses the same lexicographic minimum.  That is the step depending
on Theorem~\ref{thm:whitehead}.

\diagnostic{Distinguish ``equality is an equivalence relation'' from
``this key is complete for the intended orbit relation.''}

\subsubsection*{Answer to Exercise 6.4}

Running one implementation twice repeats the same assumptions and the same
possible bugs.  Independently written implementations with different data
flows can expose implementation-specific errors when their partitions
disagree.  Agreement is strong engineering evidence, but the claim of
completeness still rests on Whitehead's mathematical theorem; two programs
could share a mistaken theorem or the same misunderstood convention.

\diagnostic{Independent computation validates implementation, while a
proof validates the mathematical specification.}

\subsection*{Answers for Section 7: S-moves}

\subsubsection*{Answer to Exercise 7.1}

Write \(u=abcd\) and \(v=ef\), so \(r=uv\).  Then
\[
efabcd=v r v^{-1}
\]
uses a conjugator of length \(2\).  Also
\[
efabcd=u^{-1} r u
\]
uses a conjugator of length \(4\).  Both are correct; the suffix
conjugator is shorter in this case.

\diagnostic{A cyclic rotation has two obvious conjugation expressions.
Choose the shorter only for a bound, not because the other is false.}

\subsubsection*{Answer to Exercise 7.2}

There is one AC1.  Since \(\eps=-1\), there are two AC2 moves.  The AC3
count is
\[
m_1+2m_2=2+2\cdot3=8.
\]
The total is
\[
1+2+8=\boxed{11}.
\]

\diagnostic{The other relator must be conjugated back, which is why
\(m_2\) occurs twice.}

\subsubsection*{Answer to Exercise 7.3}

Starting from \((r_i,r_j)\):
\begin{align*}
(r_i,r_j)
&\xrightarrow{\AC2}(r_i,r_j^{-1})\\
&\xrightarrow{\AC3}(x r_iX,r_j^{-1})\\
&\xrightarrow{\AC3}(x r_iX,Yr_j^{-1}y)\\
&\xrightarrow{\AC1}
  (x r_iX\,Yr_j^{-1}y,\;Yr_j^{-1}y)\\
&\xrightarrow{\AC3}
  (x r_iX\,Yr_j^{-1}y,\;r_j^{-1})\\
&\xrightarrow{\AC2}
  (x r_iX\,Yr_j^{-1}y,\;r_j).
\end{align*}
The fifth line conjugates the other relator by \(y=c_2^{-1}\):
\[
y(Yr_j^{-1}y)Y=r_j^{-1}.
\]
The target retains the rotated inverse factor created before restoration.

\diagnostic{Restoring the other relator must not alter the product already
stored in the target.}

\subsubsection*{Answer to Exercise 7.4}

Using \(\phi(X)=YX\) and \(\phi(Y)=Y\):
\[
\phi(c_1)=\phi(Xy)=YXy,
\]
which is reduced and has length \(3\).  Also
\[
\phi(c_2)=\phi(xY)=xyY\longrightarrow x,
\]
which has length \(1\).  For \(\eps=+1\), the sharper bound is
\[
1+\len{\phi(c_1)}+2\len{\phi(c_2)}
=1+3+2=\boxed{6}.
\]

\diagnostic{Using only \(\Lambda(\phi)\) would give a valid but weaker
bound; reduce the complete conjugator images for the sharper value.}

\subsection*{Answers for Section 8: ACA}

\subsubsection*{Answer to Exercise 8.1}

For example,
\[
P\xrightarrow{\AC1}P_1
\xrightarrow{\phi}\phi(P_1)
\xrightarrow{\AC2}Q.
\]
AC1 and AC2 preserve AC-triviality by AC reachability, and the middle step
preserves it by Theorem~\ref{thm:main-invariance}.  Hence
\(\mathcal T(P)=\mathcal T(Q)\).  The chain does not say that \(P,Q\) are
adjacent, and the automorphism step need not itself be an ordinary AC path
between its endpoints.

\diagnostic{ACA chains preserve a property step by step; they do not
collapse every step into one AC edge.}

\subsubsection*{Answer to Exercise 8.2}

Every child generated from the chosen representative is produced by a
legal S-move and verified automorphism, so every reported edge is sound.
However, a different representative \(\phi(P)\) of the same orbit may have
different word lengths and different encoded rotations, exposing valid
children not generated from \(P\).  Those missing edges cause
incompleteness.

\diagnostic{Soundness asks whether generated edges are valid; completeness
asks whether all valid edges are generated.}

\subsubsection*{Answer to Exercise 8.3}

Removing the only connecting path leaves two reported components where the
full graph has one.  Thus the bounded search reports too many components.
Selecting one representative from each reported component may duplicate
work, but it still selects at least one representative covering every
input.  This is the safe direction.

\diagnostic{Pruning can split a true component; it cannot fuse two true
components.}

\subsubsection*{Answer to Exercise 8.4}

A raw AC-graph node is one specific canonical relator pair.  A quotient
node is the entire \(\Aut(\F)\)-orbit of such pairs, represented by its
complete Whitehead canonical key.  A quotient edge records a legal AC
move from the selected orbit representative followed by re-identification
of the child orbit.

\diagnostic{Write brackets \([P]_{\Aut}\) when speaking about quotient
nodes; this prevents confusing a presentation with its orbit.}

\subsection*{Answers for Section 9: the count pipeline}

\subsubsection*{Answer to Exercise 9.1}

\[
261-93=168
\]
uses the pure change-of-variable forest edges identifying full Aut-orbits.
Then
\[
168-42=126
\]
uses the verified ACA search edges in the canonical cap-28 sweep.  Finally,
the two probe-certified edges join two previously separate singleton pairs:
\[
126-2=124.
\]
Every subtraction is a successful union between two formerly distinct
components, not merely a listed edge.

\diagnostic{An edge lowers the component count only when its endpoints
were previously in different components.}

\subsubsection*{Answer to Exercise 9.2}

The two edges internal to existing components lower the count by zero.
The other three can lower it by at most one each, so the component count
can decrease by at most \(3\).  It may decrease by less if those three
become redundant after earlier new unions.

\diagnostic{Count successful unions, not raw edges.}

\subsubsection*{Answer to Exercise 9.3}

The bounded graph can omit a valid edge between any two reported
components.  In the extreme case, all 124 components could belong to one
true ACA component; this is compatible with the current search because
absence of an edge is not a proof.  Therefore 124 is not a lower bound.
The proved statement is that the true number is at most 124.

\diagnostic{A search-derived partition of a subgraph refines the partition
of the full graph.}

\subsubsection*{Answer to Exercise 9.4}

For each of the 20 trivialized representatives, every original input in
its certified component is proved AC-trivial.  The exact number of settled
rows is the sum of the sizes of those 20 components, which cannot be
deduced from the number 20 alone.  Nothing about the remaining 104
representatives is proved by a failed or unfinished search.

\diagnostic{Report solved components and their member counts, not simply
the number of representative runs.}

\subsubsection*{Answer to Exercise 9.5}

Repeated discovery of the same two mergers under different ceilings and
resource trajectories is evidence of reproducibility.  Exactness would
require proving that no further connecting path exists.  Because every
queue terminated on a state, memory, or time cap rather than exhaustion,
unvisited states and possible mergers remain.  Agreement cannot replace
exhaustion or a separating invariant.

\diagnostic{Repeated bounded negative results improve empirical confidence
but do not prove a global absence.}

\subsection*{Answers for Section 10: certificates}

\subsubsection*{Answer to Exercise 10.1}

Every listed edge may be valid while a separate partition-building bug
places an additional, unconnected vertex into a component.  Per-edge
verification never examines that unsupported assignment.  Rebuilding the
connected components using only verified edges checks that each reported
membership has an actual certificate chain.

\diagnostic{Edge validity and transitive-closure correctness are separate
claims and need separate tests.}

\subsubsection*{Answer to Exercise 10.2}

Let the rows of \(M\) be the exponent-sum vectors of \(r_1,r_2\).
\begin{itemize}
  \item AC1 replaces row \(i\) by row \(i\) plus row \(j\), an elementary
  unimodular row operation, so the determinant is unchanged.
  \item AC2 negates one row, so the determinant changes sign and its
  absolute value is unchanged.
  \item AC3 sends \(r_i\) to \(g r_i g^{-1}\).  The exponent contributions
  of \(g\) and \(g^{-1}\) cancel, so row \(i\) is unchanged.
\end{itemize}
Thus \(\len{\det M}\) is invariant under AC1--AC3.  A free-basis
automorphism acts on exponent sums by a matrix in
\(\mathrm{GL}(2,\mathbb Z)\), i.e.\ a unimodular column operation, so the
absolute determinant is ACA-invariant as well.

\diagnostic{Conjugation is invisible after abelianization because its
prefix and suffix have opposite exponent vectors.}

\subsubsection*{Answer to Exercise 10.3}

Constancy at \(1\) means a claimed edge joining inputs with different
values would be definitely wrong, so the invariant is useful as an
independent consistency check.  But because every input already has the
same value, it partitions the dataset into only one invariant class and
cannot prove any pair inequivalent.

\diagnostic{An invariant detects differences only when it actually takes
different values on the objects being compared.}

\subsubsection*{Answer to Exercise 10.4}

One example is to change the recorded meeting state on one side while
leaving every path step unchanged.  Literal replay should end at the
original state and fail the equality check against the corrupted claim.
Another example is to delete a forest edge but leave both endpoints in the
same reported class; the partition rebuild should reject the unsupported
membership.

\diagnostic{A good mutation targets one advertised verifier obligation and
predicts the exact failed assertion.}

\subsection*{Answers for Section 11: oral rehearsal}

\subsubsection*{Answer to Exercise 11.1}

A model two-minute response is:

\begin{quote}
Because \(F_2\) is free, specifying \(x\mapsto xy\) and \(y\mapsto y\)
defines a homomorphism.  It is invertible: the inverse sends
\(x\mapsto xY\) and \(y\mapsto y\), and both compositions fix \(x,y\).
Now take an AC path from \(P\) to \((x,y)\).  Applying the automorphism to
every relator sends AC1 and AC2 to one move of the same type.  A
generator-conjugation AC3 becomes conjugation by the finite word
\(\phi(g)\), which expands into finitely many generator conjugations.  The
transformed path ends at \((xy,y)\), a free basis, and a finite Nielsen/AC
sequence returns it to \((x,y)\).  The inverse automorphism gives the
converse.  Therefore \(P\) is AC-trivial exactly when \(\phi(P)\) is.
\end{quote}

\diagnostic{The four indispensable points are homomorphism, inverse,
transported moves, and terminal basis correction.}

\subsubsection*{Answer to Exercise 11.2}

A model response is:

\begin{quote}
Every one of the 261 inputs lies in one of 124 components connected by
verified ACA edges.  Each edge preserves AC-triviality, so one
representative per component covers every input.  The search is incomplete,
so additional valid edges may merge components; therefore 124 is an upper
bound and the true number may be smaller.  Missing edges cause redundant
work, not missing coverage.
\end{quote}

\diagnostic{Use both words ``verified'' and ``upper bound.''}

\subsubsection*{Answer to Exercise 11.3}

A rigorous response is:

\begin{quote}
I am not claiming the endpoints of every discarded run are connected by a
raw AC path.  I am using the weaker property actually needed for the
experiment: every certified AC or Aut step preserves whether the
presentation is AC-trivial.  Along a verified ACA chain, the truth value is
constant.  Therefore if I solve the retained representative, all discarded
members of that component follow.  The 137 successful union edges reduce
261 vertices to 124 certified coverage components.
\end{quote}

\diagnostic{The justification is invariance of the question, not equality
of raw AC paths.}

\subsubsection*{Answer to Exercise 11.4}

A correct response is:

\begin{quote}
No.  This analysis trivialized zero of the 261 presentations.  It grouped
the inputs into at most 124 distinct ACA problems.  The conditional result
is that if one representative from each of those 124 components is later
proved AC-trivial, all 261 presentations will be proved AC-trivial.
\end{quote}

\diagnostic{Separate classification from solution.  Never turn a component
count into a trivialization count.}

\section{Sources and reproducibility map}

\begin{longtable}{>{\raggedright\arraybackslash}p{0.29\textwidth}
                  >{\raggedright\arraybackslash}p{0.65\textwidth}}
\toprule
Purpose & Repository source\\
\midrule
\endhead
Source-paper AC1--AC3 &
\sourcefile{literature/txt/math_ml_paper_2408.15332.txt}, lines 291--302\\
Master 261-to-124 account &
\sourcefile{results/equivalence_classes/EQUIVALENCE_FINDING.md}\\
Current 124-class proof book &
\sourcefile{results/equivalence_classes/PROOFS.md}\\
Current 137-edge machine certificates &
\sourcefile{results/equivalence_classes/certificates.json}\\
Word algebra and S-move replay &
\sourcefile{experiments/equivalence_classes/lib/words.py}\\
S-move generation &
\sourcefile{experiments/equivalence_classes/lib/acmoves.py}\\
Whitehead canonicalization with witness &
\sourcefile{experiments/equivalence_classes/lib/autcanon.py}\\
Independent Whitehead implementation &
\sourcefile{experiments/analysis/whitehead.py}\\
ACA quotient search &
\sourcefile{experiments/equivalence_classes/search/aut_search.py}\\
Base certificate verifier &
\sourcefile{experiments/equivalence_classes/verify/verify_proofs.py}\\
125/124 probe verifier &
\sourcefile{experiments/equivalence_classes/verify/verify_probe_merges.py}\\
Overnight 124 artifacts &
\sourcefile{results/equivalence_classes/probe/probe_overnight_seam*.json}\\
\bottomrule
\end{longtable}

\subsection*{Classical theorem dependencies}

\begin{itemize}
  \item The universal property of free groups, proved in
  Theorem~\ref{thm:universal}.
  \item Nielsen's theorem on bases of finitely generated free groups, used
  to certify automorphisms, together with Nielsen generation
  (Theorem~\ref{thm:nielsen-generation}) to reduce the terminal basis.
  \item Whitehead's peak-reduction theorem, stated in
  Theorem~\ref{thm:whitehead}, used to make the 168 Aut-orbit count exact.
  \item The Andrews--Curtis conjecture, used only to explain why presenting
  the trivial group does not currently imply a known AC trivialization.
\end{itemize}

\end{document}
